The stepper component plays a central role in the propagation of track parameters. It is responsible for the numerical integration of the equation of motion of a charged particle. Apart from that it handles the transport of the Jacobian matrix and the covariance matrix for the track parameters. As such, the stepper is a crucial component for track reconstruction.

\section{Rework Constrained Step}

\begin{itemize}
    \color{red}
    \item what is it and why is it needed
    \begin{itemize}
        \item key points:
        \item choosing the right step length given multiple constraints
        \item overstepping resolution
    \end{itemize}
    \item what are the modifications
    \begin{itemize}
        \item key points:
        \item remove propagation direction dependence
        \item remove sign of accuracy limit
    \end{itemize}
\end{itemize}

\texttt{ConstrainedStep} is a helper class to deal with different step length requests. In the current interface of the stepper in ACTS deals with 4 different request types: \texttt{accuracy}, \texttt{navigator}, \texttt{actor}, and \texttt{user}. The \texttt{accuracy} limit is used to communicate an appropriate step length limit to achieve the requested integration accuracy target from prior \texttt{step} calls. The \texttt{navigator} limit holds the approximated distance to the next target. In case of overstepping this can also be negative. The \texttt{actor} limit holds the approximate distance to the next abort condition, for example trajectory path length limit or target surface distance. Similar to the \texttt{navigator} limit, this can also be negative. The \texttt{user} limit is a user defined limit to restrict the step size to a given maximum.

\texttt{ConstrainedStep} takes inputs from the different sources and provides a single step length. The step length is determined by the minimum of the \texttt{accuracy}, \texttt{navigator}, \texttt{actor}, and \texttt{user} limits, while overstepping resolution is always preferred over smaller absolute distances.

This class used to be quite involved and a central cause of problems for the propagation. Over multiple iterations it has been effectively rewritten to be more robust and easier to understand. One major compilcation was the handling of step direction inside \texttt{ConstrainedStep}. The propagation can be done in both directions, forward and backward. During a backward propagation we might overstep a target. To resolve this we need to reverse direction, effectively stepping forward. This was all handled inside \texttt{ConstrainedStep} and made the code fail in certain cases. In the new implementation this was simplified by removing the propagation direction sign from \texttt{ConstrainedStep}. It only deals with distances given the propagation direction as the behavior has to be symmetrical in forward and backward propagation.

Another major change was removing the sign of the \texttt{accuracy} limit. In rare cases the accuracy was set during overstep resolution which caused it get a negative sign and thus the limit was locked and could not be changed resulting in a failed propagation. The new implementation decides the sign for the accuracy limit based on the \texttt{navigator} and \texttt{actor} limits.

\section{Adaptive step size}

\begin{itemize}
    \color{red}
    \item what is it and why is it needed
    \item what are the modifications
\end{itemize}

Adaptive step size is a technique to improve the performance of the numerical integration given a certain accuracy target. The idea is to adjust the step size based on the local curvature of the trajectory. In regions of high curvature a smaller step size is used to resolve the trajectory better. In regions of low curvature a larger step size is used to speed up the propagation and to limit the accumulation of machine precision errors.

This is done by computing an error estimate after each step of the numerical integration and comparing it to the requested accuracy target. If the error estimate is larger than the requested accuracy target the step size is reduced. If the error estimate is smaller than the requested accuracy target the step size is increased.

\section{EigenStepper}
\label{sc:eigenstepper}

The \texttt{EigenStepper} is a long standing stepper implementation in ACTS with the goal of being readable, maintainable and fast. It uses the C++ Eigen library for linear algebra operations which is the de facto standard. The math was once derived by hand and then transcribed into C++ code.

While some of the math operations are quite readable and maintainable, the extension mechanism diminishes these qualities. The idea behind the extension mechanism is to have one stepper implementation for different kinds of material and interation environments. There is a \texttt{DefaultExtension} which handles the integration in a vakuum and a \texttt{DenseExtension} which handles the integration in a dense material. Another layer inbetween, \texttt{StepperExtensionList}, allowed to switch between any number of extensions depending on their priority. This complexity was removed during the work on this thesis and the \texttt{EigenStepper} was simplified to deal with only a single extension, while the dense extension will automatically fall back to the default extension if no volume material is present.

\section{SympyStepper}

\begin{itemize}
    \color{red}
    \item motivation
    \item comparison of EigenStepper, AtlasStepper, SympyStepper approach
    \item how is this done
\end{itemize}

The \texttt{SympyStepper} is a novel stepper implementation in ACTS, conceptualized and implemented as part of this thesis. It is based on generated code from the symbolic math library \texttt{SymPy}. The motivation for this stepper was to improve CPU performance while keeping or improving the readable representation of the mathematics.

The \texttt{EigenStepper} uses the C++ Eigen library for linear algebra operations which results in readable and hopefully fast code, as the library is sort of a de facto standard for linear algebra in C++ and supports vectorization. There are still a few weeknesses in terms of readability and maintainability with this implementation which is discussed in section \ref{sc:eigenstepper}. However, it was observed that the \texttt{AtlasStepper} is significantly faster than the \texttt{EigenStepper}. The \texttt{AtlasStepper} is based on the Athena numerical implementation STEP and does not rely on any external libraries. The math operations are decomposed by hand and optimized for performance. This indeed results in fast code but is also practically unreadable and impossible to maintain. For example introducing a time parameter in the \texttt{AtlasStepper} would be significantly more involved than in the \texttt{EigenStepper}. It was attempted to improve the performance of the \texttt{EigenStepper} but turned out to be quite difficult and potentially reducing again readability and maintainability.

The \texttt{SympyStepper} can be thought of as the midpoint between the \texttt{EigenStepper} and the \texttt{AtlasStepper}. The stepper math is written in Python using the \texttt{SymPy} library. This allows for high level linear algebra operations while also deriving the linear error propagation from the Runge-Kutta-Nystrom integration. Apart from that we can use common subexpression elimination to simplify the math upfront and reduce the number of operations. Another big advantage is that we can encode known values and symmetries which will automatically result in simplified terms and fewer operations. Known values can be zeros in Jacobians we know from the beginning. We can also make use of the symmetry of the covariance matrix to avoid duplicate work. These kinds of simplifications are not possible with Eigen. Finally, we can generate C++ code from the \texttt{SymPy} expressions which can be used in the ACTS stepper implementation. The generated code is not meant to be readable but fast, can be thought of as binary code produced by a compiler and could be generated by a build system during the compilation of ACTS.

\section{The cost of the time parameter}

\begin{itemize}
    \color{red}
    \item time is not switchable
\end{itemize}

As mentioned before, the track parametrization is not customizable in ACTS. As such, the time parameter will always be present, even if it is not needed. This is the case for the ATLAS ITk detector which does not provide timing information. Since track reconstruction is a performance critical task, and we do additional operations on the time parameter, we need to be aware of the cost of the time parameter.

Given the situation in ACTS, it is not easy to estimate this cost and any shortcuts might skew the results. For this reason the time parameter was consistently removed from ACTS in a separate branch and the performance was measured. Using the \texttt{EigenStepper} the performance was measured to be 10\% faster without the time parameter. This is a significant improvement and shows that the time parameter is indeed costly. However, after the implementation of the \texttt{SympyStepper} the impact was measured to be below 1\%. This suggests that the \texttt{SympyStepper} is indeed able to optimize the math operations to a high degree as the time parameter adds mostly zeros in derivatives for the linear error propagation.
